% !PRINTED 219
\weberpage{219}

Die Wurzeln von $f$ entsprechen den sechs Octaëderecken oder
den sechs Würfelflächen; die Wurzeln von $W$ den acht Octaëderflächen
oder Würfelecken, und die Wurzeln von $K$ sowohl beim Octaëder
als beim Würfel den 12 Kanten (vgl. \S.\,48).

Wenn man die Formen $W$, $K$ so darstellt:
\[
W = (x_1^4 + x_2^4)^2 + 12\,x_1^4x_2^4, \quad K = (x_1^4 + x_2^4)^3 - 36\,x_1^4x_2^4(x_1^4 + x_2^4),
\]
so lässt sich leicht die zwischen den drei Functionen $f$, $W$, $K$
bestehende identische Relation ableiten:
\weq{15}{W^3 - K^2 = 108\,f^4.}
Die Octaëdergruppe enthält, wie man sieht, die Tetraëdergruppe
und entsteht aus ihr durch Hinzunahme der einen Substitution
 $\Theta(x) = ix$. Die Grundformen des Octaëders sind also
zugleich invariante Formen des Tetraëders, und in der That
stimmen die Formen $f$ des Tetraëders und Octaëders mit einander
 genau überein, und es ist $W = \Phi_1\Phi_2$,
\[
K = \tfrac{1}{2}(\Phi_1^3 + \Phi_2^3), \quad (\S.\,56).
\]
Mit der Determinante 1 geschrieben, lauten die beiden erzeugenden
 Substitutionen der Octaëdergruppe so:
\[
\Theta = \begin{pmatrix} \sqrt{i}, & 0 \\ 0, & \dfrac{1}{\sqrt{i}} \end{pmatrix}, \quad
\chi = \begin{pmatrix} \dfrac{1}{\sqrt{2i}}, & \dfrac{1}{\sqrt{2i}} \\[6pt] \dfrac{1}{i\sqrt{2i}}, & \dfrac{-1}{i\sqrt{2i}} \end{pmatrix},
\]
worin $\sqrt{2i} = 1 + i$ zu setzen ist.

Durch die Substitution $\Theta$ \gesperrt{ändert} nun die Form $f(x_1, x_2)$,
wie die Formel (3) unmittelbar zeigt, ihr \gesperrt{Vorzeichen}. Durch
Anwendung von $\chi$ werden die linearen Factoren von $f$ folgendermaassen
 verändert:
\[
\begin{gathered}
x_1, \quad x_2, \quad x_1 + x_2, \quad x_1 - x_2, \quad x_1 + ix_2, \quad x_1 - ix_2 \\
\frac{x_1 + x_2}{\sqrt{2i}}, \quad \frac{x_1 - x_2}{i\sqrt{2i}}, \quad -i(x_1 + ix_2), \quad (x_1 - ix_2), \quad \frac{2x_1}{\sqrt{2i}}, \quad \frac{2x_2}{\sqrt{2i}},
\end{gathered}
\]
und daraus geht hervor, dass $f(x_1, x_2)$ durch Anwendung der
Substitution $\chi$ \gesperrt{ungeändert bleibt}, und dadurch sind die
Aenderungen der Form $f$ durch alle anderen Octaëdersubstitutionen
 zugleich mit bestimmt.

Die \gesperrt{Hesse}'sche Covariante $W$ von $f$ bleibt ungeändert,
wenn $f$ in $-f$ verwandelt wird, und folglich bleibt $W$ bei den
%CONT>>
